% Requires \usepackage{listings} in the main document preamble.
\lstdefinestyle{pythonTC}{
    language=Python,
    basicstyle=\ttfamily\scriptsize,
    numbers=none,
    frame=single,
    breaklines=true,
    columns=fullflexible,
    keepspaces=true,
    showstringspaces=false,
    aboveskip=4pt,
    belowskip=4pt
}

\section{Computational Implementation}
\label{app:computational-implementation}

The core numerical algorithms were implemented explicitly with Python lists,
loops, arithmetic operations, and the standard library. NumPy is used for the
input-array construction specified in the assignment and to generate plotting
coordinates; Matplotlib is used only to render the figure. The solution does
not call high-level library solvers such as \texttt{numpy.linalg.qr},
\texttt{numpy.linalg.lstsq}, \texttt{numpy.linalg.solve}, or
\texttt{numpy.polyfit}. These abridged excerpts rely on helper routines also
defined in the source modules; the complete executable files include
validation, diagnostics, command-line interaction, and report formatting.

\subsection{QR Factorization and Least-Squares Solution}

The inner products, norms, vector updates, and matrix operations used by the
solver are implemented with loops and summations. The reduced QR factorization
uses Modified Gram--Schmidt: each column is successively orthogonalized
against the previously computed columns of $Q$. The projection coefficients
are stored in the upper-triangular matrix $R$.

\begin{lstlisting}[style=pythonTC,
    caption={Core operations and Modified Gram--Schmidt factorization from
    \texttt{linear\_algebra.py}.},
    label={lst:mgs-core}]
import math
import sys


def vector_dot(u, v):
    return sum(ui * vi for ui, vi in zip(u, v))


def vector_norm(u):
    return math.sqrt(sum(ui * ui for ui in u))


def qr_mgs(A, tol=None):
    A = to_float_matrix(A)
    m, n = matrix_shape(A)
    scale = frobenius_norm(A)
    if tol is None:
        tol = max(m, n) * sys.float_info.epsilon * scale

    Q = [[0.0 for _ in range(n)] for _ in range(m)]
    R = [[0.0 for _ in range(n)] for _ in range(n)]
    q_cols = []
    full_rank = True

    for j in range(n):
        v = get_column(A, j)
        for i in range(j):
            qi = q_cols[i]
            R[i][j] = vector_dot(qi, v)
            v = [v[k] - R[i][j] * qi[k] for k in range(m)]

        R[j][j] = vector_norm(v)
        if R[j][j] <= tol:
            full_rank = False
            qj = [0.0] * m
        else:
            qj = [v[k] / R[j][j] for k in range(m)]
        q_cols.append(qj)
        for k in range(m):
            Q[k][j] = qj[k]

    info = {
        "is_full_rank": full_rank,
        "tolerance": tol,
        "diagonal_entries": [R[i][i] for i in range(n)],
    }
    return Q, R, info
\end{lstlisting}

The rank threshold is scale-dependent,
\[
    \mathrm{tol}=\max(m,n)\,\epsilon_{\mathrm{mach}}\,\|A\|_F.
\]
If a diagonal entry of $R$ does not exceed this threshold, the
implementation flags possible numerical rank deficiency and avoids
normalizing by a negligible value. Modified Gram--Schmidt without column
pivoting is not a universally reliable rank-revealing method for severely
ill-conditioned matrices.

For full-column-rank matrices, $A=QR$ reduces the least-squares problem to
$R\mathbf{x}=Q^T\mathbf{b}$. The program forms $Q^T\mathbf{b}$ with its own
matrix--vector routine and solves the triangular system by backward
substitution.

\begin{lstlisting}[style=pythonTC,
    caption={Back substitution and the QR-based least-squares solve.},
    label={lst:backsub}]
def back_substitution(R, d, tol):
    n = len(R)
    x = [0.0] * n
    for i in range(n - 1, -1, -1):
        if abs(R[i][i]) <= tol:
            raise RankDeficiencyError("Numerically singular system")
        s = sum(R[i][j] * x[j] for j in range(i + 1, n))
        x[i] = (d[i] - s) / R[i][i]
    return x


def solve_least_squares_qr(A, b, tol=None):
    Q, R, info = qr_mgs(A, tol)
    if not info["is_full_rank"]:
        raise RankDeficiencyError("Possible rank deficiency")
    d = matrix_vector_mul(matrix_transpose(Q), b)
    x = back_substitution(R, d, info["tolerance"])
    return x, Q, R, info
\end{lstlisting}

\subsection{Data Construction and Quadratic Fit}

For Problem~1, the matrix and right-hand-side vector retain the expressions
and NumPy array construction supplied in the assignment. That prescribed
constructor is not reproduced here. For Problem~2, the design matrix is
assembled as rows $(x_i^2,x_i,1)$ using Python list operations; the same QR
solver is reused, and the MSE is calculated directly from the residuals.

\begin{lstlisting}[style=pythonTC,
    caption={Problem 2: design matrix, fitted values, and explicit MSE.},
    label={lst:tc3-data-mse}]
# x and y are generated according to the assignment.
X = [[float(xi)**2, float(xi), 1.0] for xi in x]
theta, Q, R, info = solve_least_squares_qr(X, y)
a, b_coef, c = theta
y_hat = [a*xi**2 + b_coef*xi + c for xi in x]
mse = sum((yi - pi)**2 for yi, pi in zip(y, y_hat)) / len(y)
\end{lstlisting}

\subsection{Numerical Verification and Visualization}

The verification metrics are computed using the same explicit matrix and
vector routines. For each matrix $B$ and right-hand side $\mathbf{c}$, the
implementation evaluates
\begin{align}
    \eta_{\mathrm{QR}} &= \frac{\|B-QR\|_F}{\|B\|_F}, &
    \eta_Q &= \|Q^TQ-I\|_F,\\
    \rho &= \|B\mathbf{z}^*-\mathbf{c}\|_2, &
    \eta_{\mathrm{LS}} &= \|B^T(B\mathbf{z}^*-\mathbf{c})\|_2.
\end{align}
Here, $\eta_{\mathrm{QR}}$ and $\eta_Q$ assess reconstruction and
orthogonality, whereas $\eta_{\mathrm{LS}}$ checks the least-squares
optimality condition. The residual norm and MSE measure fit to the data,
not numerical error in the solver. Matplotlib is used only to render the
curve and observations; it does not estimate the polynomial coefficients.
